% ED26-BM-OHTR2012-4: actual Markman chapter4, printed36--51.
\chapter{Cognitive Knowledge Representation}
\label{sec:p19-cognitive-knowledge-representation}

A cognitive representation is not only a stored object. A theory must
specify what it represents, how it receives its interpretation and
what processes can use it. Markman's chapter distinguishes mental
spaces, feature representations and structured representations,
including semantic networks and relational structures. It then asks
how knowledge content, bodily states and environmental context alter
that picture, and answers antirepresentational objections through
representational pluralism
\cite[chap.~4, pp.~36--51]{Markman2012CognitiveRepresentation}.

The central question is how a format makes particular reasoning
operations possible, easy, difficult or unavailable. A successful
formal encoding does not establish that people use it; a measured
response does not by itself identify the encoding or algorithm that
produced it. The source's cognitive arguments, empirical surveys and
mathematical dependencies retain these different standings below.

\section{Represented information, representing states and interpretation}
\label{sec:p19-ckr-definition}

\begin{hbdefinition}{def:p19-ckr-four-components}{Four components of a functioning representation}
Following Markman's reconstruction of Palmer and Pylyshyn, a
representation comprises a \emph{representing world}, the domain
serving as the representation; a \emph{represented world}, the
information or domain represented; \emph{representing relations}
connecting them and supplying interpretation; and \emph{processes}
that use the representing states for some function. The represented
domain can be external information or other internal representations;
it is not restricted to a currently visible object.

A collection of marks is not yet evidence that its user possesses
every fact an observer can read into it. The interpretation and
usable operations must belong to the proposed system, not merely
to the theorist viewing its diagram. Markman's usage is broad enough
to include nonmental representations. His narrower mental usage
locates the representing world inside the head, whereas photographs,
writing and tools can be external resources. The source's later
embodied/situated discussion addresses their coordination without
silently changing this initial mental/external distinction.
\end{hbdefinition}

Figure~4.1 uses digit symbols for numerical quantities and separately
specifies order and increment/decrement procedures. Its represented
world includes real numbers, while its simple discrete code does
not retain all their information. Neither understanding the drawn
digits nor recognizing their intended association supplies a missing
operation to the modeled system. Likewise a photograph can represent
a scene through a lawful image-formation relationship, but the
viewer's perceptual/conceptual procedures do the work of interpreting
it. Markman's point is the completeness of the representation-plus-
process explanation, not a theorem that every digital encoding must
lose every task-relevant property
\cite[pp.~36--38 and fig.~4.1]{Markman2012CognitiveRepresentation}.

The historical motivation is the computational view of mind that
gained influence in the cognitive revolution. A data structure and
its algorithm can be described without specifying all hardware
details. Markman places representation and algorithm at Marr's
algorithmic level, distinguishes the abstract function at the
computational level and the physical realization at the implementation
level, and recognizes constraints from the latter. These are
explanatory levels, not findings that all brains implement identical
programs or that implementation can never matter.

\subsection{Grounding and conceptual role}

The thermostat illustrates grounding: changes in temperature cause
differential expansion of two metals, changing a strip's curvature;
that curvature can operate a switch. The lawful relation and the
device's use of it support a temperature interpretation. In perception,
some internal representations are grounded through others: retinal
transduction precedes further processing such as edge detection.
A chain of representing relations can connect internal states with
external conditions \cite[p.~38]{Markman2012CognitiveRepresentation}.

But reference to a currently present stimulus cannot explain every
mental content. Markman contrasts seeing his dog with imagining it
while it is elsewhere, and mentions nonexistent entities such as
unicorns. Conceptual-role semantics offers another source of meaning:
an element's role among other representations. The dictionary analogy
then limits an exclusively internal explanation. Looking up unfamiliar
words only produces more unfamiliar words; some links to nonlinguistic
information are needed for understanding on this account. The source
therefore combines grounding and conceptual relations rather than
declaring either sufficient for every content. It cites Fodor,
Stich and Warfield and Searle as philosophical dependencies; their
full theories are not replaced by the analogy.

\section{Dimensions that do not define one universal format}
\label{sec:p19-ckr-dimensions}

The chapter's three families are \emph{spatial}, \emph{featural}
and \emph{structured}. Other distinctions crosscut them. A state
can be continuous or composed of discrete symbols; its elements
can be independent or explicitly linked; it can track momentary
conditions or persist over time; and its interpretation can be
analog or symbolic. These dimensions must not be collapsed into
a single distinction between pictures and words
\cite[p.~39]{Markman2012CognitiveRepresentation}.

In the source's analog-clock example, angular displacement represents
elapsed time within the clock's bounded cycle. A 90-degree interval
has half the represented duration of a 180-degree interval under
the stipulated scale. That conclusion uses an actual structural
correspondence. Wraparound, which hand is being read and the time
span assigned to a revolution remain interpretation conditions.
Digit symbols instead need explicit procedures for corresponding
numerical relations; their shapes do not themselves implement
addition or comparison.

Analogical structure is not the same as a conscious visual image.
Nor is distributed activation a fourth mutually exclusive alternative
to all three families. The source cites interactive-activation and
symbolic-connectionist implementations alongside symbolic models.
A pattern spread over units, a labelled graph and an imagined
layout need separately specified contents, bindings and operations.
An implementation distributed across units does not by itself
determine whether it retains the role of each object in a relation.

\section{Spatial representation: proximity and its constraints}
\label{sec:p19-ckr-spatial}

The representing world can be a multidimensional space. Objects
occupy points or regions; coordinates locate them. A physical map
illustrates two-dimensional representation, but a semantic space
need not have only three dimensions or describe geographical
locations. Interpreting a coordinate or a distance as a semantic
property requires its own representing relation.

\begin{hbdefinition}{def:p19-ckr-distance}{The source's distance, with its mathematical conditions}
Let \(N\geq1\) be an integer, let \(x,y\in\mathbb R^N\)
have coordinates \(x_i,y_i\), and let \(r\geq1\) be a finite
real exponent. Here absolute-value bars denote the magnitude of a
real coordinate difference, not cardinality. The Minkowski distance is
\[
 d(x,y)=\left(\sum_{i=1}^N \lvert x_i-y_i\rvert^r\right)^{1/r}.
\]
The city-block case is \(r=1\); Euclidean straight-line distance
is \(r=2\). Absolute values and the exponent domain are essential
for the general metric claim. They repair equation~1's unrestricted
printed expression rather than reproduce its missing conditions.
\end{hbdefinition}

\begin{hbtheorem}{thm:p19-ckr-metric}{proved=annex:p19.ckr-metric}{The metric guarantee and its exponent boundary}
Under Definition~\ref{def:p19-ckr-distance}, distances are nonnegative,
zero exactly for equal points, symmetric and satisfy the triangle
inequality. The same coordinate formula with \(0<r<1\) need
not satisfy the triangle inequality. A metric representation with
fixed coordinates cannot realize dissimilarities that actually violate
that inequality; an unspecified qualitative similarity description
alone does not establish such a violation.
\end{hbtheorem}

Multidimensional scaling, MDS, seeks a spatial approximation to
supplied proximity information. Markman's European-cities example
illustrates recovery of a geographical configuration from suitable
distances, not a guarantee that every input yields an exact or
unique map. In Figure~4.2, the two panels place bird and mammal
concepts in semantic spaces: close conceptual neighbors and distant
items have different positions. The chapter attributes similarity
ratings and sentence-verification contrasts to Rips, Shoben and
Smith 1973. Its robin/duck illustration reports faster verification
for the robin's category membership. Its accompanying phrase
\emph{inversely related to distance} is not consistent with that
illustration and is not promoted here to an established quantitative
latency law. The chapter gives no participant denominator, task
assignment details, latency estimates or uncertainty interval for
that comparison \cite[pp.~40--41]{Markman2012CognitiveRepresentation}.

\begin{hbtheorem}{thm:p19-ckr-pairs}{proved=annex:p19.ckr-pairs}{The source's pair-rating count}
For a nonnegative integer \(X\) and \(X\) distinct labelled items,
one judgment for each
unordered pair of different items requires exactly
\(X(X-1)/2\) judgments. With 15 items this is 105, and with
20 items it is 190. The count does not apply unchanged if
directional judgments, repetitions or a selected subset are used.
\end{hbtheorem}

This accounts for equation~2 and the source's concern about elicitation
effort. It does not prove a universal maximum of 20 represented
concepts. Corpus-based approaches discussed through Burgess and Lund
and Landauer and Dumais instead derive high-dimensional semantic
spaces from word co-occurrence. Markman's sentence with an unspecified
actor jumping upstairs illustrates contextual constraints on what
word can occupy a position. Those methods and MDS have different
inputs; the chapter points outside for their full algorithms. No
particular fitting objective, optimization guarantee, neural mechanism
or contemporary embedding model is supplied by the shared word
\emph{space} \cite[p.~41]{Markman2012CognitiveRepresentation}.

Spatial formats have two proposed advantages: data-driven construction
and simple proximity computations. Their limitation is that distances
or difference vectors do not automatically tell the modeled system
which named property is shared. An observer's interpretation of an
axis cannot be assigned to the system without a mechanism.

The moon/ball/lamp example makes the issue salient. Roundness makes
moon and ball similar; brightness makes moon and lamp similar;
different commonalities are salient across pairs. The source relates
such findings to Tversky and Gati and to triangle-inequality failures.
To establish a failure, however, a proposed numerical dissimilarity
must exceed the sum of the other two. For an explicitly editorial
example, distances 1, 1 and 3 cannot describe one fixed metric
triangle. Qualitative labels such as moderate similarity and no
similarity do not alone determine those numbers. Context-sensitive
attention or modified spaces, discussed through Krumhansl and
Nosofsky, can alter predictions but also add process complexity.
Neither simplicity nor a more flexible fit uniquely identifies cognition
\cite[pp.~41--42]{Markman2012CognitiveRepresentation}.

\section{Feature representation: discrete properties and missing bindings}
\label{sec:p19-ckr-features}

A feature representation contains discrete items referring to
properties. The source's speech example contrasts /b/ and /d/
by place of articulation, and /b/ and /p/ by voicing. A proposed
speech-perception system must link such features to acoustic
information and combine them to identify a phoneme. Finding
a linguistically useful contrast does not itself identify an
invariant acoustic detector. The chapter explicitly notes the
difficulty of that mapping, citing Jakobsen, Fant and Halle
and Blumstein and Stevens
\cite[p.~42]{Markman2012CognitiveRepresentation}.

Tversky's surveyed contrast account compares feature sets through
intersection and set difference: shared features increase similarity,
distinctive features reduce it. Markman reports feature-list and
similarity-judgment relationships, but supplies no sample, task
denominator, effect coefficient or interval. Its formal set operations
are specified; the human weighting and feature-elicitation claims
remain a source-survey boundary, not a new participant result.

Features can be augmented with centrality or typicality and with
degrees of belief. The source's bird/dolphin example distinguishes
category-defining properties from properties merely characteristic
of a category. It contrasts an overall comparison with a further
diagnostic comparison when classification is uncertain. Its wording
about focusing on characteristic properties conflicts with the
subsequent defining/core-property explanation; the latter distinction
is retained without inventing a verified two-stage algorithm.
The Smith, Shoben and Rips 1974 model is an external dependency.

Likewise, certainty factors do not automatically explain why a
property is believed or why it matters to a category. The Lenat/Guha
and Shafer pointers identify examples rather than one shared
probability interpretation or combination rule. To represent grounds
and causal relations explicitly, a system needs more than an
unconnected list of properties
\cite[pp.~42--43]{Markman2012CognitiveRepresentation}.

\section{Structured representation: argument roles and usable relations}
\label{sec:p19-ckr-structured}

Relations bind their arguments in specific roles. Poodles being a
kind of dog does not assert the converse; someone taller than
another person occupies a different role from the person being
compared. Markman's source notation uses question-prefixed slots:
\(\operatorname{kind\!-\!of}(?x,?y)\),
\(\operatorname{red}(?x)\) and
\(\operatorname{above}(?x,?y)\). A one-argument attribute identifies
its scope, while a two-argument relation identifies both roles.
These are cognitive representation schemata. They are not by
themselves the exact first-order grammar or an evaluated logical
calculus \cite[p.~43]{Markman2012CognitiveRepresentation}.

\begin{hbexample}{def:p19-ckr-scenes}{One task applied to the source's three families}
Figure~4.3 contrasts a striped circle above a shaded square with
a striped square above a shaded circle. Both can receive the same
flat feature list: above, circle, square, round, angled, striped,
shaded. Ask which shape is striped and above the other. The
following is an editorial matched-task comparison using the source
scenes, not a newly conducted experiment.

In a labelled spatial layout, the system can inspect the upper
region and its pattern/shape, provided those properties and their
spatial locations are actually preserved and a region-inspection
procedure is available. A generic semantic-space distance alone
does not supply that layout or inspection mechanism.

In the specified flat feature bag, both scenes are encoded alike.
Comparing their intersection or the presence of striped cannot
recover which object carries it. Object-indexed or compound features
could repair this loss, but would change the encoding rather than
make the missing binding appear from the unchanged list.

In the bound relational representation, the first scene assigns
striped to the circle and shaded to the square, and supplies the
ordered above relation; the second reverses the shape roles.
A query can follow the attribute's object argument and the upper
role. Success still requires a process that follows those connections.
Displaying a graph for an observer is not proof that an implemented
system executes the query or that a person uses the graph.
\end{hbexample}

\begin{hbdefinition}{def:p19-ckr-decoding}{An explicit task-information comparison}
For this local formalization let \(\mathcal S\) be a nonempty
finite set of scenes, \(e:\mathcal S\to\mathcal R\) an encoding
into representing states, and \(q:\mathcal S\to\mathcal A\)
the task's correct answer, with \(\mathcal A\) nonempty.
A decoder for the task is a function on the actually used codes
\(e(\mathcal S)\) taking values in \(\mathcal A\).
These are editorial mathematical objects for assessing retained
information, not measured neural states or a cognitive algorithm.
\end{hbdefinition}

\begin{hbtheorem}{thm:p19-ckr-decoding}{proved=annex:p19.ckr-decoding}{Exactly when an encoding permits this task's answers}
A decoder \(g:e(\mathcal S)\to\mathcal A\) satisfies
\(q=g\circ e\) exactly when scenes with the same code always
have the same answer. Thus the flat bag in
Example~\ref{def:p19-ckr-scenes} admits no decoder giving the
correct shape for both source scenes. This is a task-relative
information-loss guarantee, not a result that every feature
representation is inadequate or that a successful decoder is
efficient or psychologically implemented.
\end{hbtheorem}

\subsection{Semantic networks: marker passing is not activation spreading}

Figure~4.4 links vampires, hematologists, blood, water, clothing and
superordinate categories. Directed labelled links can encode vampires
drinking blood and hematologists studying it. In marker passing,
markers retain their originating nodes and propagate through
outgoing links. An intersection at blood allows tracing the two
paths and retrieving the different labelled relationships. Under
the described synchronized procedure, finding an intersection at
a given step depends on available directed paths, not on an
unexplained impression that the objects resemble one another
\cite[p.~44]{Markman2012CognitiveRepresentation}.

Spreading activation instead changes node accessibility. Activation
can travel both ways along links, and link strengths can differ.
The source's doctor/nurse priming example is attributed to Meyer
and Schvaneveldt; the chapter reports no latency estimate, participant
population, trial structure or uncertainty. A priming contrast is
not by itself evidence for this unique spreading mechanism.
In some Anderson-style models, a fixed activation quantity is
divided over a node's outgoing links, its fan. A larger fan then
allocates less to each neighbor under that rule. Fan division is
one process assumption, not a property of every graph or every
network implementation. Marker identity/path tracing, graded
accessibility, interactive activation and explanatory coherence
have different operations even when diagrams all show linked nodes
\cite[pp.~44--45]{Markman2012CognitiveRepresentation}.

\subsection{Relational systems, scripts and analogy}

Verbs and prepositions bind actors, objects and locations into
relations. Scripts and schemas add order and causal organization
to familiar events. Markman's restaurant case separates the expected
arrival of a menu from knowing why the menu is useful: if none
arrives, the diner can seek another source of the restaurant's
offerings, such as a wall board. Merely knowing an event sequence
does not specify the purpose-sensitive response to its interruption
\cite[p.~45]{Markman2012CognitiveRepresentation}.

The source's solar-system/atom comparison illustrates analogy
between structurally related but superficially dissimilar domains.
The better-known base is the solar system; the target is the atom.
Its note records that base is also called source. Matching the
revolve-around relation aligns planet with electron in its first
argument, and sun with nucleus in its second: parallel connectivity
constrains the object correspondences. A one-to-one constraint
allows each base object at most one target match. Systematicity
favors connected systems, including higher-order causal relations,
over isolated similarities
\cite[pp.~45--46 and note~1]{Markman2012CognitiveRepresentation}.

The correctly named atom schemata include
\[
 \operatorname{revolve\!-\!around}(\mathrm{electron},\mathrm{nucleus}),
\]
\[
 \operatorname{greater}(\operatorname{mass}(\mathrm{nucleus}),
                        \operatorname{mass}(\mathrm{electron})).
\]
A corresponding source-domain causal schema is
\[
\begin{aligned}
 \operatorname{cause}\bigl(&
 \operatorname{greater}(\operatorname{mass}(\mathrm{sun}),
                         \operatorname{mass}(\mathrm{planet})),\\
 &\operatorname{revolve\!-\!around}(\mathrm{planet},\mathrm{sun})\bigr).
\end{aligned}
\]
This repairs the source's misnaming of the first domain and missing
parenthesis. The higher-order schema takes relational contents as
arguments; it is not silently declared well formed in P2's
first-order language. Nor is its familiar classical atom analogy
a modern physical law. Transferring the causal relation to the
target yields a candidate analogical inference that need not be
true. A mapping constraint supplies neither truth of the base's
causal premise nor warranted transfer of all its properties.

Markman associates structure-sensitive processing with more
constraints and greater resource demands than elementary distance
or feature comparison. He points to the Structure-Mapping Engine
and other analogy models for algorithmic detail and complexity.
The chapter gives no universal computational bound for every
possible relational encoding. The full model and its chosen
search/control procedure remain external dependencies; a relation
graph alone is not the missing algorithm.

\section{Content, bodily states and situated activity}
\label{sec:p19-ckr-embodiment}

Format does not exhaust knowledge. Markman contrasts the relatively
content-independent computational tradition with developmental
research on what children know and when they can use it. He
also points to expert/novice differences and pragmatic reasoning
schemas. His observation that many cognitive studies use college
undergraduates is a source-level methodological assessment, not
a sampled proportion of the literature or a finding about every
student. Age, domain knowledge, expertise and task interpretation
cannot be ignored when transferring a representational account
\cite[p.~47]{Markman2012CognitiveRepresentation}.

Embodied approaches challenge a clean separation in which perception
only supplies input, cognition computes abstractly, and motor control
only executes output. In the Gibson comparison, vision serves
action and object affordances: a chair can be assessed for sitting,
standing or lifting. This differs from merely reconstructing a
veridical three-dimensional scene. The source's action-oriented
proposal is a theoretical alternative, not proof that seeing a
chair always entails a specific conscious calculation.

Markman surveys three kinds of evidence. Wu and Barsalou's property-
generation example contrasts a lawn with a rolled-up lawn, which
makes roots more salient. The proposed account uses perceptual
simulation; verbal associations and task demands are competing
mechanisms that the example alone does not eliminate. Hill-slope
reports are described as varying with a heavy backpack, citing
Proffitt, Creem and Zosh. A further friend/social-support contrast
is described without a separate identified study citation in this
chapter. Finally the chapter points to gesture research for links
between abstract thought and motor information. It provides no
participant denominators, recruitment populations, allocation details,
effect estimates or uncertainty intervals for those examples.
They remain precise attributed survey claims, not independently
qualified causal effects or evidence that bodily load uniquely
identifies the representation used
\cite[pp.~47--48]{Markman2012CognitiveRepresentation}.

Situated cognition adds task context and external tools. Lists
can support fallible memory; Hutchins's ship-navigation example
emphasizes how team activity is organized by its instruments.
External information, bodily states and internal representations
are different proposed contributors, not interchangeable names.
The claim that a task is supported by tools does not prove that
no internal representation is used. Nor does interaction with the
world mean the representation is a faithful model of every relevant
feature. The question is how the parts coordinate their actual
information and operations.

\section{Antirepresentationalism, its insight and the pluralist reply}
\label{sec:p19-ckr-pluralism}

The source identifies Port and Van Gelder, Spivey and Thelen and
Smith with challenges to the computational/representational tradition.
One target is a plan-then-execute picture: construct a detailed
complete plan first and pass it to a separate executor. Real
behavior must respond during execution to changing conditions.
Markman's route-to-a-friend illustration changes a route when
construction closes a road. A static finished plan would not
by itself specify that response
\cite[pp.~48--49]{Markman2012CognitiveRepresentation}.

Reactive robotics demonstrates that some sophisticated behavior,
such as hallway traversal with obstacle avoidance, can be obtained
without enduring world states supporting elaborate advance plans.
The chapter cites Beer, Brooks and Pfeifer and Scheier. In its
own broad four-component definition, however, these systems can
still possess minimal representations. Success without a particular
kind of stored plan is not the elimination of every possible
representing state. The disputed conclusion and the accepted
architectural criticism must be separated.

Markman accepts the importance of online coordination but rejects
extending its success into a universal no-representation account.
Kuipers's spatial-semantic hierarchy supplies his constructive
example: a changing-environment level, a large-scale planning
level and intermediate representations connecting plans with
motor programs. This is a source-described robot architecture
and reported implementation, not an acquired codebase or a
newly verified benchmark. Its full interfaces and performance
evidence are external to the selected chapter.

The pluralist conclusion follows as a research proposal: distinct
formats have different advantages, and a coordinated system may
need several. Fast changing-state control and enduring relational
abstraction need not use one narrow format; parsimony does not
require ignoring task-defining differences. Multiple encodings
alone are also insufficient: their interfaces and control must
actually support robust behavior. The source leaves three explicit
future questions: how to study knowledge content more fully,
how bodily/environmental representations interact with abstract
ones, and how multiple representations coordinate within one model
\cite[p.~49]{Markman2012CognitiveRepresentation}.

\section{Formal interfaces, evidence and consultation}
\label{sec:p19-ckr-consultation}

The first-order structure definition in \cref{sec:fol-semantics-01}
supplies a nonempty domain and interpretations of the nonlogical symbols;
\cref{sec:fol-semantics-02} gives the separate satisfaction clauses.
The grammar distinguishes terms from formulas.
That is not a model of how people acquire grounding or retrieve
a semantic relation. P9's Definition~\ref{def:p09-alc} supplies
an explicit description-logic interpretation of concepts and roles;
decidability or completeness of a formal system is not evidence
that a person uses its procedures. Cognitive predicate diagrams
also require an actual syntax and interpretation before a formal
guarantee can be transferred to them.

Johnson-Laird's Chapter~\ref{sec:p19-rationality-2-3} concerns
construction and search of mental possibilities. Its iconicity
does not require a conscious picture, and its formal live-model
benchmark is distinct from a person's incomplete representation.
Here Markman compares formats more broadly. A task solved by
a possibility model does not determine every representation in
the broader cognitive architecture.

\begin{hbmethod}{def:p19-ckr-use}{Determine what the encoding and process actually supply}
Specify the represented information, representing states,
interpretation and operations. Identify which task answer or
inference is required. Do not import the observer's knowledge
into the modeled system merely because a diagram is intelligible.

Test whether different relevant cases receive the same code.
The striped-shape task exposes a lost binding in the flat bag.
If adopting a spatial model, specify the metric, fitted inputs
and mapping from proximity to the observed judgment; qualitative
similarity is not already a distance scale. For a relational
model, specify roles, matching constraints, control and treatment
of candidate inferences. Surface resemblance, path intersection
and truth-preserving deduction are different operations.

Distinguish a representation's adequacy from evidence of its use.
Compare rival encodings and algorithms that could produce the
same response; investigate their discriminating task predictions.
For evidence, retain population, task, design, operational measure,
uncertainty and generalization conditions as in
Method~\ref{def:p19-evidence-record}. A better fit or a successful
robot does not uniquely identify human cognition.

For bodily or environmental contributions, state what information
changes and how it reaches the process. Distinguish an absent
long-term plan from absent representation, and distinguish a
multiple-format vocabulary from an implemented interface. Preserve
source-specific empirical boundaries before recommending transfer
to another task, age group, occupation or neural substrate.
\end{hbmethod}

The successful geometric-scene query uses retained object/property
bindings and an appropriate inspection or graph procedure. The
failed flat-bag query lacks sufficient information, independent
of processing speed. Conversely an information-sufficient encoding
may remain computationally unusable under a deadline, or may
not be the encoding a participant actually constructs. These
are different failure cases. The annex's metric, pair-count and
decoder results establish only their stated mathematical claims;
the source's broader cognitive conclusions retain their argued
and precisely surveyed standing.

\index{knowledge representation!cognitive formats}
\index{representation!grounding and conceptual role}
\index{representation!spatial}
\index{representation!features and bindings}
\index{representation!structured relations}
\index{representation!embodied and situated}
\index{representational pluralism}
